
\documentclass{article}

% -----------------------------
% Packages
% -----------------------------
\usepackage{fancyhdr}
\usepackage{extramarks}
\usepackage{amsmath}
\usepackage{amsthm}
\usepackage{amsfonts}
\usepackage{tikz}
\usepackage[plain]{algorithm}
\usepackage{algpseudocode}
\usepackage{amssymb}
\usepackage{enumitem}

\usetikzlibrary{automata,positioning}

% -----------------------------
% Basic Document Settings
% -----------------------------
\topmargin=-0.45in
\evensidemargin=0in
\oddsidemargin=0in
\textwidth=6.5in
\textheight=9.0in
\headsep=0.25in

\linespread{1.1}

\pagestyle{fancy}
\lhead{\hmwkAuthorName}
\chead{\hmwkClass: \hmwkTitle}
\rhead{}
\lfoot{\lastxmark}
\cfoot{\thepage}

\renewcommand\headrulewidth{0.4pt}
\renewcommand\footrulewidth{0.4pt}

\setlength\parindent{0pt}

% -----------------------------
% Problem Section Helpers
% -----------------------------
\newcommand{\enterProblemHeader}[1]{
  \nobreak\extramarks{}{Problem #1 continued on next page\ldots}\nobreak{}
  \nobreak\extramarks{Problem #1 (continued)}{Problem #1 continued on next page\ldots}\nobreak{}
}

\newcommand{\exitProblemHeader}[1]{
  \nobreak\extramarks{Problem #1 (continued)}{Problem #1 continued on next page\ldots}\nobreak{}
  \nobreak\extramarks{Problem #1}{}\nobreak{}
}

\setcounter{secnumdepth}{0}
\newcounter{partCounter}
\newcounter{homeworkProblemCounter}
\setcounter{homeworkProblemCounter}{1}
\nobreak\extramarks{Problem \arabic{homeworkProblemCounter}}{}\nobreak{}

% -----------------------------
% Theorem env (optional)
% -----------------------------
\newtheorem*{theorem}{Theorem}

% -----------------------------
% Homework Problem Environment
%   Usage:
%   \begin{homeworkProblem}            % auto-number
%     % Your answer here
%   \end{homeworkProblem}
%
%   \begin{homeworkProblem}[1.2.3]     % label problem "1.2.3" (e.g. textbook numbering)
%     % Your answer here
%   \end{homeworkProblem}
% -----------------------------
\newenvironment{homeworkProblem}[1][]{
  \def\hwProblemLabel{#1}
  \ifx\hwProblemLabel\empty
    \edef\hwProblemLabel{\arabic{homeworkProblemCounter}}
    \stepcounter{homeworkProblemCounter}
  \fi
  \section{Problem \hwProblemLabel}
  \setcounter{partCounter}{1}
  \enterProblemHeader{\hwProblemLabel}
}{
  \exitProblemHeader{\hwProblemLabel}
}

% -----------------------------
% Homework Details (EDIT THESE)
% -----------------------------
\newcommand{\hmwkTitle}{Homework \#7}
\newcommand{\hmwkClass}{Math 3320}
\newcommand{\hmwkAuthorName}{\textbf{Joseph Quinn}}

% -----------------------------
% Title
% -----------------------------
\title{
  \vspace{2in}
  \textmd{\textbf{\hmwkClass:\ \hmwkTitle}}\\
  \vspace{3in}
}
\author{\hmwkAuthorName}
\date{}

% -----------------------------
% Convenience Macros (optional)
% -----------------------------
\renewcommand{\part}[1]{\textbf{\large Part \Alph{partCounter}}\stepcounter{partCounter}\\}

% Algorithms
\newcommand{\alg}[1]{\textsc{\bfseries \footnotesize #1}}

% Calculus
\newcommand{\deriv}[1]{\frac{\mathrm{d}}{\mathrm{d}x} (#1)}
\newcommand{\pderiv}[2]{\frac{\partial}{\partial #1} (#2)}
\newcommand{\dx}{\mathrm{d}x}

% Statistics
\newcommand{\E}{\mathrm{E}}
\newcommand{\Var}{\mathrm{Var}}
\newcommand{\Cov}{\mathrm{Cov}}
\newcommand{\Bias}{\mathrm{Bias}}

% Proof step (for aligned derivations)
\newcommand{\step}[2]{& #1 & & \text{#2} \\}

% Solution header (optional)
\newcommand{\solution}{\textbf{\large Solution}}

% -----------------------------
% Document
% -----------------------------
\begin{document}

\maketitle
\pagebreak



\begin{homeworkProblem}[10.2.5]
	First thing we want to do is notice the ciphertext contains VHV twice,
	at position 0 and position 16. This difference of 16 suggests that
	the key length divides 16, giving us the possible key lengths
	\[
		\{1,2,4,8,16\}.
	\]

	The next clue we get is that the ciphertext AE represents plaintext AN.
	This means the key value used to encode the A is
	\[
		\begin{aligned}
			0+x & \equiv 0 \pmod{26} \\
			x   & \equiv 0 \pmod{26} \\
			x   & =0
		\end{aligned}
	\]

	And for E
	\[
		\begin{aligned}
			13+x & \equiv 4 \pmod{26}  \\
			x    & \equiv -9 \pmod{26} \\
			x    & \equiv 17 \pmod{26} \\
			x    & =17
		\end{aligned}
	\]

	This tells us the key uses the letters A and R at positions 9 and 10.
	This immediately rules out the key having length 1, since a key of
	length 1 would use the same letter at every position.

	The next clue is that VHV represents a common 3 letter word,
	which I'm going to assume is THE.

	Using the same process, for the first letter V representing T,
	we get
	\[
		\begin{aligned}
			19+x & \equiv 21 \pmod{26} \\
			x    & \equiv 2 \pmod{26}  \\
			x    & =2
		\end{aligned}
	\]

	This gives us the key letter C.

	For the second letter H representing H,
	\[
		\begin{aligned}
			7+x & \equiv 7 \pmod{26} \\
			x   & \equiv 0 \pmod{26} \\
			x   & =0
		\end{aligned}
	\]

	This gives us the key letter A.

	And for the third letter V representing E,
	\[
		\begin{aligned}
			4+x & \equiv 21 \pmod{26} \\
			x   & \equiv 17 \pmod{26} \\
			x   & =17
		\end{aligned}
	\]

	This gives us the key letter R.

	So assuming VHV represents THE, we now know the first
	three letters of the repeating key are CAR.

	This also rules out key length 2, since positions 0 and 2
	would need to use the same key letter. However, we found
	that position 0 uses C and position 2 uses R.

	Now we need to determine if the key is length 4, 8, or 16 and to do this we can use frequency
	analysis. Since every letter encrypted by the same position
	in the key uses the same Caesar shift, we can separate the
	ciphertext into groups based on the possible key length.

	We then compare the letter frequencies in each group with
	English letter frequencies to find the most likely shifts.

	I used the website `\texttt{https://www.dcode.fr/vigenere-cipher}' to test this and found the most likely
	key length was 4, which produces the key CARD.

\end{homeworkProblem}
\begin{homeworkProblem}[10.3.3]

	We are given a Feistel network with the following information
	\begin{itemize}
		\item $n=4$
		\item 2 rounds
		\item
		      \[
			      \begin{aligned}
				      f_0(x_1,x_2,x_3,x_4) & := (x_2 \oplus x_4,\ 1,\ x_1x_2,\ 1 \oplus x_3), \\
				      f_1(x_1,x_2,x_3,x_4) & := (1,\ x_1 \oplus x_3,\ x_4,\ x_2).
			      \end{aligned}
		      \]
		\item $k=(0,1)$, which tells us the function to use at each round
	\end{itemize}

	So following the standard encryption route, to encrypt $m=m_0m_1$

	\[
		\begin{aligned}
			m_2 & = m_0 \oplus f_0(m_1) \\
			m_3 & = m_1 \oplus f_1(m_2)
		\end{aligned}
	\]

	we are given the ciphertext $c=m_3m_2=10101011$, so
	\[
		m_3=1010,\qquad m_2=1011
	\]


	To decrypt, we reverse the rounds:

	\[
		\begin{aligned}
			m_1 & = m_3 \oplus f_1(m_2) \\
			m_0 & = m_2 \oplus f_0(m_1)
		\end{aligned}
	\]


	so first we recover $m_1$:
	\[
		\begin{aligned}
			m_1 & = m_3 \oplus f_1(m_2)          \\
			    & = 1010 \oplus f_1(1011)        \\
			    & = 1010 \oplus (1,1\oplus1,1,0) \\
			    & = 1010 \oplus 1010             \\
			    & = 0000
		\end{aligned}
	\]

	and then we recover $m_0$:
	\[
		\begin{aligned}
			m_0 & = m_2 \oplus f_0(m_1)                       \\
			    & = 1011 \oplus f_0(0000)                     \\
			    & = 1011 \oplus (0\oplus0,1,0\cdot0,1\oplus0) \\
			    & = 1011 \oplus 0101                          \\
			    & = 1110
		\end{aligned}
	\]

	Therefore, the original message is
	\[
		m=m_0m_1=11100000
	\]
\end{homeworkProblem}

\begin{homeworkProblem}[11.1.11(a)]

	\begin{proof}
		From the division algorithm we know for any integers $a$ and $b$
		with $b>0$, there exist integers $q$ and $r$ such that
		\[
			a=qb+r,\qquad 0\leq r<b.
		\]

		And in the Euclidean algorithm, we repeatedly divide the
		previous remainder by the current remainder. Letting $r_0=n$ and
		$r_1=m$, this gives
		\[
			\begin{aligned}
				n   & = q_1m+r_2   \\
				m   & = q_2r_2+r_3 \\
				r_2 & = q_3r_3+r_4 \\
				r_3 & = q_4r_4+r_5
			\end{aligned}
		\]

		So any given remainder \( r_i \) we can formulate as
		\[
			r_i=q_{i+1}r_{i+1}+r_{i+2}
		\]

		Obviously we know that the remainders are decreasing so,
		$r_i>r_{i+1}$.

		Since $r_{i+1}<r_i$, the quotient satisfies $q_{i+1}\geq 1$, so
		\[
			\begin{aligned}
				r_i    & = q_{i+1}r_{i+1}+r_{i+2} \\
				r_i    & \geq r_{i+1}+r_{i+2}
			\end{aligned}
		\]

		So knowing the fact that \( r_{i+1} > r_{i+2} \)
		we can add $r_{i+2}$ to both sides giving
		\[
			r_{i+1}+r_{i+2}>2r_{i+2}.
		\]

		From the inequality we formed above we know \( r_i \) is greater than this making
		\[
			r_i\geq r_{i+1}+r_{i+2}>2r_{i+2}.
		\]

		so we can just cut out the middle man leaving
		\[
			r_i>2r_{i+2}.
		\]

		and with division we get
		\[
			r_{i+2}<\frac{r_i}{2}
		\]
	\end{proof}
\end{homeworkProblem}

\begin{homeworkProblem}[11.1.12(b)]

	\begin{proof}
		\[
			\begin{aligned}
				([a]+[b])[c]
				 & =[a+b][c]       \\
				 & =[(a+b)c]       \\
				 & =[ac+bc]        \\
				 & =[ac]+[bc]      \\
				 & =[a][c]+[b][c].
			\end{aligned}
		\]

		Therefore, multiplication distributes over addition

	\end{proof}
\end{homeworkProblem}

\begin{homeworkProblem}[11.1.13]
	$\gcd(105,180)$

	\[
		\begin{aligned}
			180 & = 1(105)+75 \\
			105 & = 1(75)+30  \\
			75  & = 2(30)+15  \\
			30  & = 2(15)+0
		\end{aligned}
	\]

	Since the last nonzero remainder is $15$, we have

	\[
		\gcd(105,180)=15
	\]

	We now work backward to express the GCD as a linear combination of $105$ and $180$.

	\[
		15=75-2(30).
	\]

	sub in $30=105-75$

	\[
		\begin{aligned}
			15 & = 75-2(105-75)    \\
			   & = 75-2(105)+2(75) \\
			   & = 3(75)-2(105).
		\end{aligned}
	\]

	sub in $75=180-105$

	\[
		\begin{aligned}
			15 & = 3(180-105)-2(105)    \\
			   & = 3(180)-3(105)-2(105) \\
			   & = 3(180)-5(105).
		\end{aligned}
	\]

	And now we are done as it's in terms of 105 and 180,
	\[
		105(-5)+180(3)=\gcd(105,180)=15.
	\]
\end{homeworkProblem}

\begin{homeworkProblem}[11.1.15]
	\begin{enumerate}[label=(\alph*)]
		\item Let $n=6$, $a=4$, and $b=3$.

		      We have
		      \[
			      ab=4\cdot 3=12
		      \]
		      Since $6\mid 12$, we have $n\mid ab$. However,
		      \[
			      6\nmid 4 \quad \text{and} \quad 6\nmid 3.
		      \]
		      Therefore, $n\mid ab$ does not imply $n\mid a$ or $n\mid b$.

		\item Let $a=7$ and $p=4$.

		      Since $(7,4)=1$, the first condition is satisfied. However,
		      \[
			      a^{p-1}=7^{4-1}=7^3=343.
		      \]
		      Since
		      \[
			      343\equiv 3\pmod 4,
		      \]
		      we have
		      \[
			      7^3\not\equiv 1\pmod 4.
		      \]
		      Therefore, $(a,p)=1$ does not imply
		      $a^{p-1}\equiv 1\pmod p$.

		\item Let $a=6$, $b=4$, and $c=2$.

		      We have
		      \[
			      (ab,c)=(24,2)=2.
		      \]
		      However,
		      \[
			      (a,c)(b,c)=(6,2)(4,2)=2\cdot 2=4.
		      \]
		      Since
		      \[
			      2\neq 4,
		      \]
		      we conclude that
		      \[
			      (ab,c)\neq (a,c)(b,c).
		      \]
		      Therefore, the statement is false in general.
	\end{enumerate}

\end{homeworkProblem}
\begin{homeworkProblem}[A1]

	\begin{proof}
		Suppose $a$ and $b$ are both inverses of $c$, so that
		$c*a=e$ and $b*c=e$. Then
		\[
			\begin{aligned}
				b & = b*e     \\
				  & = b*(c*a) \\
				  & = (b*c)*a \\
				  & = e*a     \\
				  & = a
			\end{aligned}
		\]

		Therefore, $a=b$, so inverses are unique.

	\end{proof}
\end{homeworkProblem}

\end{document}
